\documentclass[tikz,border=4pt]{standalone}
\usepackage{ctex}
\usepackage{amsmath,amssymb}
\usetikzlibrary{calc,arrows.meta,decorations.markings,patterns}
\begin{document}
\begin{tikzpicture}[
  >=Stealth,
  font=\small,
  ccw/.style={->, very thick, decoration={markings,
      mark=at position 0.13 with {\arrow{>}},
      mark=at position 0.38 with {\arrow{>}},
      mark=at position 0.63 with {\arrow{>}},
      mark=at position 0.88 with {\arrow{>}}},
    postaction={decorate}}
]
\draw[->] (-1.5,0) -- (5.8,0) node[right] {$x$};
\draw[->] (0,-1.5) -- (0,5) node[above] {$y$};

% 区域 D（用椭圆，简洁）
\fill[blue!10] (2.5, 2.0) ellipse [x radius=1.8, y radius=1.4];

% 内部小点示意"区域积分"
\foreach \xx in {1.3, 2.0, 2.7, 3.4} {
  \foreach \yy in {1.3, 2.0, 2.7} {
    \fill[blue!30] (\xx, \yy) circle (1.5pt);
  }
}

% 边界 L 正向（逆时针）
\draw[ccw, blue!70!black] (2.5, 2.0) ellipse [x radius=1.8, y radius=1.4];

% 标签
\node[blue!50!black, font=\large] at (2.5, 2.0) {$D$};
\node[blue!70!black, anchor=west] at (4.5, 1.0) {$L=\partial D$};
\node[blue!70!black, font=\footnotesize, anchor=west] at (4.5, 0.55) {(正向 = 逆时针)};

% "区域始终在左侧" 提示（在椭圆右上沿切向画左指法向）
\draw[->, thick, green!60!black] (4.05, 2.7) -- ++(-0.5, 0.4);
\node[green!50!black, font=\footnotesize, anchor=west] at (3.4, 3.5) {区域在左侧};

% 公式（左上空白处，远离椭圆）
\node[font=\small, draw=gray!60, fill=gray!10, fill opacity=0.9,
      text opacity=1, rounded corners=2pt, align=left, anchor=west]
  at (-1.4, 4.5)
  {\textbf{Green 公式}\\[2pt]
   $\displaystyle\oint_L P\,dx+Q\,dy=\iint_D\!\!\left(\frac{\partial Q}{\partial x}-\frac{\partial P}{\partial y}\right)dA$\\[2pt]
   \footnotesize 边界环量 = 内部"旋度"积分};

% 底部说明
\node[font=\footnotesize, align=center] at (2.5, -1.1)
  {统一模式：$\displaystyle\int_{\partial D}\!\!\bullet \;=\; \iint_D\!\bigl[\text{对 }\bullet\text{ 求"内部导数"}\bigr]$};
\end{tikzpicture}
\end{document}
